%% ma: L10-B13-0106
%% lop: 10
%% chuong: 5
%% bai: 13
%% dang: 10.13.13
%% loai: DS
%% muc: [NB, TH, TH, VD]
%% dap_an: "ĐSĐĐ"
%% nguon: "chuyen-de-cac-so-dac-trung-cua-mau-so-lieu-khong-ghep-nhom-toan-10-chuong-trinh-moi.pdf (GV. Phan Nhật Linh) (P03), Bài 02 Dạng 2 Phần II Câu 10"
%% kien_thuc: "Lớp 10 Bài 13 (số trung bình, trung vị, tứ phân vị, mốt)"
%% trang_thai: chinh-thuc
%% ghi_chu: "Tự duyệt theo yêu cầu giáo viên 22/10/2026."
\begin{ex}
Khối lượng cơ thể lúc trưởng thành của $11$ con chim được ghi lại ở bảng sau (đơn vị: gam):
\begin{center}
\begin{tabular}{|c|c|c|c|c|c|c|c|c|c|c|}\hline
$165$ & $150$ & $155$ & $165$ & $170$ & $165$ & $150$ & $155$ & $170$ & $165$ & $160$\\\hline
\end{tabular}
\end{center}
Xét tính đúng sai của các khẳng định sau (các kết quả làm tròn đến hàng đơn vị):
\choiceTF{\True Mốt của mẫu số liệu là $165$ gam}{Tứ phân vị dưới của mẫu số liệu là $160$}{\True Khối lượng trung bình của $11$ con chim là $161$ gam}{\True Mẫu số liệu tập trung với mật độ cao ở bên phải $Q_2$ và mật độ thấp ở bên trái $Q_2$}
\loigiai{Sắp xếp: $150$; $150$; $155$; $155$; $160$; $165$; $165$; $165$; $165$; $170$; $170$.

a) Đúng: $165$ xuất hiện $4$ lần.

b) Sai: $Q_2=165$; nửa dưới $150$; $150$; $155$; $155$; $160$ nên $Q_1=155$.

c) Đúng: $\overline x=\dfrac{1770}{11}\approx160{,}9\approx161$.

d) Đúng: $Q_1=155$, $Q_2=165$, $Q_3=165$: khoảng từ $Q_2$ đến $Q_3$ bằng $0$, từ $Q_1$ đến $Q_2$ bằng $10$.}
\end{ex}
